\scanpage{259}
\mspage{209}

localiseurs $W_S,W_{S'}$. La \deleted{\unsure{façon}} notion de $W$-fibration sert donc en premier lieu à dégager la condition la plus évidente sur une flèche $f:S'\to S$ dans $M$, pour que \deleted{\unsure{elle}} le foncteur \unsure{changement|chgt} de base \deleted{\unsure{induit}} soit compatible avec les localiseurs $W_S,W_{S'}$ (c'est \uline{ça} la définition), donc induise $(*)$ sur les \unsure{localisées}.

Le ``\uline{problème des \deleted{\unsure{types}} fibres homotopiques}'' est justement de définir \unsure{(x)|(*)} pour \uline{toute} \deleted{\unsure{classe}}\inserted{flèche} dans $M$, pas seulement pour les $W$-fibrations, de telle façon que les $\mathrm{Ho}^{!}_W(S)$ forment, pour $S$ variable, les fibres d'une catégorie fibrée remarquable sur $M$. Quand $M$ admet un objet final \inserted{$e$}, \unsure{de} sorte que l'on peut regarder les flèches
\[
p:e\unsure{\longrightarrow|\text{---}}S\qquad\text{($p$ comme ``point'')}
\]
dans $M$ comme des ``points'' de $S$, la ``fibre $W$-homotopique'' d'un objet quelconque $X/S$ de $\deleted{\unsure{\mathrm{Cat}|C}}\unsure{M}_{/S}$ en ledit point sera par définition l'élément de
\[
\mathrm{Ho}_W=\mathrm{Ho}_W(M)\simeq\mathrm{Ho}^{!}_W(e)
\]
défini par
\[
p^*(\mathrm{ho}^{!}_{W/S}(X/S)),
\]
% archive sheet number: 259
